We presented a continuous hidden Markov model (CHMM) with Gaussian emissions trained via the Baum-Welch algorithm for generating synthetic financial time series, together with two cross-asset dependence extensions for multi-asset synthesis.
The central univariate finding is that the CHMM reproduces all three canonical stylized facts of equity returns (heavy-tailed distributions, negligible linear autocorrelation, and persistent volatility clustering) at moderate state resolution ($K = 6$--$13$) \emph{without} requiring jump augmentation, discretization, or any mechanism beyond the standard HMM framework.

This result directly challenges the influential conclusion of \citet{ryden1998stylized} that Gaussian HMMs cannot reproduce the slow decay of the autocorrelation function of absolute returns.
We demonstrated that this limitation is an artifact of insufficient state resolution: \citet{ryden1998stylized} tested only $K = 2$--$3$ states, where the geometric sojourn-time distribution cannot approximate slow ACF decay.
At $K \geq 3$ with quantile-based initialization, the Baum-Welch algorithm learns a transition matrix with sufficient diagonal dominance in high-volatility states to sustain persistence, achieving ACF-MAE of 0.046--0.049 (comparable to GARCH(1,1) at 0.047) through a sum-of-exponentials approximation that effectively mimics the flexible sojourn-time distributions of hidden semi-Markov models \cite{bulla2006stylized}.

The elimination of the jump mechanism relative to the discrete framework of \citet{alswaidan2026hybrid} removes two hyperparameters ($\epsilon$ for jump probability, $\lambda$ for mean jump duration) and the associated grid search, reducing computational cost and model complexity.
The joint optimization of emission parameters and transition probabilities via Baum-Welch allows the model to learn tail-state persistence directly from data, making the external jump forcing unnecessary.

The four-model benchmark confirmed that no single generator dominates all seven evaluation metrics, but the CHMM achieves the most balanced performance: 94.4\% KS pass rate (vs.\ GARCH's 13\%), ACF-MAE of 0.048 (vs.\ GARCH's 0.047), Wasserstein-1 of 0.11 (vs.\ Laplace's 0.12), and 100\% quantile coverage.
Cross-asset univariate generalization to NVDA (97.4\% IS KS), JNJ (99.8\%), and JPM (98.6\%) confirmed adaptability to diverse equity risk profiles, and out-of-sample evaluation on held-out 2025 data showed robust generalization with KS pass rates above 91\% at all state resolutions.

The cross-asset extensions then showed that the CHMM marginals compose cleanly with standard cross-sectional dependence models.
On a six-asset universe, both the Gaussian and Student-t copulas substantially outperformed the Single Index Model factor baseline: copulas achieved a median in-sample KS pass rate of 98.0\% versus 77.0\% for SIM, and an off-diagonal correlation MAE of 0.030 (Gaussian) and 0.026 (Student-t) versus 0.079 for SIM.
The Student-t copula's profile MLE selected $\nu^* = 6$ degrees of freedom, and the resulting tail-dependent copula produced the best correlation reproduction overall.
These findings replicate the ordering reported by \citet{alswaidan2026hybrid} with discrete HMM marginals, suggesting that the copula advantage is intrinsic to the rank-reordering construction and transfers across marginal families.

Extension of the framework to VIX log returns demonstrated that the CHMM captures the distinct statistical properties of volatility indices (positive skewness, mean reversion, and extreme kurtosis) and provides a natural bridge to option pricing through regime-switching geometric Brownian motion parameterized by CHMM-decoded volatility levels.
A detailed comparison of the CHMM-implied volatility surface against a parametric stochastic-volatility benchmark is the subject of companion work led by our group, and intentionally lies outside the scope of this paper.

Several directions for future work follow from these results:
\begin{enumerate}
    \item \textit{Heavy-tailed emissions.}
    Replacing Gaussian emissions with Student-$t$ or generalized hyperbolic distributions within the Baum-Welch framework would address the kurtosis gap (7.71 observed vs.\ 3.7--5.1 simulated) while retaining joint optimization of emissions and transitions.

    \item \textit{Time-varying transition dynamics.}
    Estimating the transition matrix via rolling windows or Bayesian online learning would relax the stationarity assumption and potentially improve OoS performance for assets with non-stationary regime structures.

    \item \textit{Full-universe multi-asset generation via vine copulas.}
    Scaling the copula construction from six assets to the 424-asset universe of \citet{alswaidan2026hybrid} requires vine or factor copulas for tractability and is a natural next step.

    \item \textit{Model selection.}
    Applying information-theoretic criteria (AIC, BIC) or cross-validation to select $K$ would provide principled guidance on state resolution, complementing the sensitivity analysis presented here.
\end{enumerate}

\paragraph{Data and Code Availability.}
The simulation scripts, training and test data, and all code to reproduce the results in this paper are available under an MIT license from the paper repository: \url{https://github.com/altashly1/CHMM-paper}.
The core continuous HMM implementation, the SIM module, and the Gaussian / Student-t copula modules are available in the development repository: \url{https://github.com/altashly1/HMM-withJumps-WIP}.

\paragraph{Conflict of Interest.}
The authors declare no competing interests.

\paragraph{Author Contributions.}
J.V.\ directed the study.
A.A.\ developed the continuous model, implemented the Baum-Welch algorithm, implemented the Single Index Model and copula-based cross-asset extensions, conducted the empirical analysis, and generated all figures and tables.
All authors edited and reviewed the final manuscript.
